%% =====================================================================
%%  卒業論文テンプレート（日本語版）
%%  法政大学 理工学部 経営システム工学科  知能システム研究室
%% ---------------------------------------------------------------------
%%  ・コンパイル：LuaLaTeX（Overleaf では Menu → Settings → Compiler）
%%  ・書き換えるもの：このファイルの「論文情報」，chapters/ の各章，
%%                    references.bib（参考文献），figures/（図）
%%  ・hosei-thesis.cls と latexmkrc は編集しないこと
%%  ・詳しい使い方は README.md を参照
%% =====================================================================

%  クラスオプション
%    ja / en            本文の言語（日本語 / 英語）
%    bachelor / master  卒業論文 / 修士論文
%    review             指導教員に見せるときに付ける
%                       （行番号・\TODO・\rev の色付けを表示．提出する版では外す）
\documentclass[ja,bachelor]{hosei-thesis}

%% ---------------------------------------------------------------------
%%  論文情報（表紙に印刷される）
%% ---------------------------------------------------------------------
%  日本語は単語の途中でも改行されるので，題目が長いときは \\ で改行位置を指定する
\titleJA{○○を用いた△△問題の解法}                     % 論文題目（和文）
\titleEN{Solving the YY Problem Using XX}            % 論文題目（英文）
\author{法政 太郎}             % 氏名
\authorkana{ホウセイ タロウ}    % フリガナ
\studentid{22X0000}           % 学生証番号
\supervisor{劉 子昂}           % 指導教員
\fiscalyear{2026}             % 提出する年度に変更し，定稿時に固定する（4月〜翌年3月）

%% ---------------------------------------------------------------------
%%  追加のパッケージや自作コマンド（必要なときだけ）
%%  よく使うもの（図・表・数式・アルゴリズム・ソースコード）は読み込み済み
%% ---------------------------------------------------------------------
% \usepackage{...}
% \newcommand{\vect}[1]{\boldsymbol{#1}}

\begin{document}

\maketitle                         % 表紙

\frontmatter                       % 前付け（ページ番号は i, ii, ...）
%% =====================================================================
%%  要旨
%%  研究の背景・課題・目的・方法・結果・結論を 1 ページ（800〜1200 字程度）にまとめる．
%%  要旨だけを読んでも研究の内容がわかるように書く．
%%  図・表・数式・文献の引用は使わない．
%%  下の例の「%% 背景」などは構成の目安を示すコメント（PDF には出ない）．
%%  「（背景）」のような見出しを本文に書く必要はない．
%%  修士論文の場合，研究科に別途提出する「論文要旨」（2〜8 頁，Hoppii の様式）は，
%%  この要旨とは別の書類である．
%% =====================================================================
\begin{abstractJA}
\begin{writingnote}
要旨は本文を書き終えてから最後に書き，本文にない内容は入れない．
欠かせないのは目的と主な結果である．背景と課題は 1〜2 文にとどめ，方法と結果に字数を割く．
結果は「有効性を確認した」のような抽象的な表現で済ませず，どのような条件で何がわかったかを具体的に書く（主要な数値を示すとよい）．
最後に，結果から言えることと，提案手法の適用条件や限界を述べる．
略語は初出で正式な名称を示す．用語と数値は本文，特に「結論」の章と一致させる．
\end{writingnote}

%% 背景
近年，……の分野では，○○が重要になっている．
%% 課題
しかし，従来の方法には……という課題がある．

%% 目的・方法
そこで本研究では，……を目的として，○○を用いた△△問題の解法を提案する．
提案手法では，……により……する．

%% 結果・結論
数値実験により，……の条件で……という結果が得られた．
以上より，提案手法の適用条件として……，限界として……が明らかになった．
\end{abstractJA}

%% 英文の要旨（Abstract）も付ける場合は，次の行頭の % を外す．
% \begin{abstractEN}
% Write the abstract in English here.
% \end{abstractEN}
          % 要旨
\tableofcontents                   % 目次
% \listoffigures                   % 図目次（必要なら）
% \listoftables                    % 表目次（必要なら）

\mainmatter                        % 本文（ページ番号は 1, 2, ...）
%% =====================================================================
%%  第1章 はじめに
%%  書くこと：背景 → 課題 → 目的 → 提案の概要 → 貢献 → 本論文の構成
%%  この章だけで「なぜこの研究が必要で，何をしたのか」がわかるようにする．
%% =====================================================================
\chapter{はじめに}
\label{chap:introduction}

\section{研究の背景}
\label{sec:background}

\begin{writingnote}
研究の背景を述べる．社会的・学術的な文脈から書き始め，本研究が扱う問題へと話題を絞っていく．
事実や主張には根拠を示し，必要に応じて文献を引用する．
題目に含まれる用語などの主要な用語は，初めて使うところで簡単に説明する（厳密な定義は「準備」の章で述べる）．

%% 段落を変えるときは空行を 1 行入れる．\\ で改行しないこと．
次の段落では，主な既存研究の成果に簡単に触れたうえで，まだ解決されていない課題を明らかにする．
既存研究の詳しい整理は「関連研究」の章で行う．
読者が「その課題を解決する必要がある」と納得できるように，具体例や数値を挙げて説明するとよい．

%% 論文の構成や文章の書き方は，次の参考書が役に立つ（引用の例も兼ねている）．
論文の構成や文章の書き方については，文献~\cite{kinoshita1981rikakei,sakai2015kore}が参考になる．
\end{writingnote}

近年，……において○○が課題となっている．
既存の方法では……であるため，……を検討する必要がある．

\section{研究の目的}
\label{sec:purpose}

\begin{writingnote}
本研究の目的を 1〜2 文で明確に述べる．
目的は前の節で述べた課題に対応させ，「結論」の章で達成できたかどうかを判断できるように具体的に書く．
続いて，目的を達成するための提案手法の概要と，実験で何を確かめたかを簡潔に述べる．
優位性を前提にせず，実験で確認できた性能，適用条件，限界を貢献として整理する．
この節は，ほかの章を書き終えてから見直し，目的と貢献が「提案手法」と「数値実験」の章で示した内容と一致しているかを確かめる．
\end{writingnote}

本研究の目的は，……を明らかにすることである．
本研究の貢献は次のとおりである．
\begin{itemize}
  \item ○○問題を△△として定式化した．
  \item ○○を用いた△△問題の解法を提案した．
  \item 数値実験により，提案手法の性能，適用条件および限界を明らかにした．
\end{itemize}

\section{本論文の構成}
\label{sec:organization}

%% 章・節・図・表・式は \label{名前} で名前を付け，\cref{名前} で参照する．
%% 「第2章」「3.1節」「図4.1」「式(3.1)」などの語と番号は自動で付く（手で書かない）．
本論文の構成は以下のとおりである．
\cref{chap:preliminaries}では，本研究で扱う問題の定義と，提案手法の基礎となる既存の解法について述べる．
\cref{chap:related-work}では，関連研究を整理し，本研究の位置づけを明らかにする．
\cref{chap:method}では，提案手法について述べる．
\cref{chap:experiments}では，数値実験の結果を示し，考察する．
\cref{chap:conclusion}では，本論文をまとめ，今後の課題を述べる．
    % はじめに
%% =====================================================================
%%  第2章 準備
%%  書くこと：本研究の前提となる既存の知識（扱う問題の定義，用いる手法の基礎）
%%  ここは自分の提案ではない．教科書や論文の内容を，出典を明記して簡潔にまとめる．
%%  自分の提案は「提案手法」の章に書き，既存の内容とはっきり分ける．
%%  以下では例として，0-1 ナップサック問題とその貪欲法を説明している．
%%  自分の研究に合わせて書き換える．
%% =====================================================================
\chapter{準備}
\label{chap:preliminaries}

本章では，本研究で扱う問題を定義し，提案手法の基礎となる既存の解法について述べる．
\begin{writingnote}
この章には，提案手法と数値実験を理解するために必要な既存の知識だけを書く．教科書のように網羅する必要はない．
定義や既存の解法は，原典（元の論文や標準的な教科書）を引用して説明する．
文献によって定義が異なる用語は，本論文で用いる定義を明記し，以後はその意味だけで使う．
記号は分野の標準的な教科書に合わせ，一度定義した記号を別の意味で使わない．主な記号は\cref{chap:notation}のように一覧にするとよい．
\end{writingnote}

\section{対象とする問題}
\label{sec:problem}

%% 問題は，言葉による説明 → 数式による定式化 → 記号の説明 の順に書くとわかりやすい．
%% 複数行の数式は align 環境で書き，& を置いた位置で揃える．&& の後には条件を書く．
%% 連続した式は \crefrange{最初}{最後} で「式(2.1)〜式(2.3)」のように参照できる．
例として，0-1 ナップサック問題~\cite{kellerer2004knapsack}を考える．
0-1 ナップサック問題は，価値と重さが与えられた $n$ 個の品物の中から，
重さの合計が容量を超えないように品物を選び，価値の合計を最大化する問題であり，
\crefrange{eq:knapsack-objective}{eq:knapsack-binary}のように定式化できる．
\begin{align}
  \text{maximize}   \quad & \sum_{i=1}^{n} v_i x_i \label{eq:knapsack-objective} \\
  \text{subject to} \quad & \sum_{i=1}^{n} w_i x_i \le W \label{eq:knapsack-capacity} \\
                          & x_i \in \{0, 1\} && i = 1, \dots, n \label{eq:knapsack-binary}
\end{align}
ここで，$v_i$ と $w_i$ はそれぞれ品物 $i$ の価値と重さ，$W$ は容量であり，
$x_i$ は品物 $i$ を選ぶとき 1，選ばないとき 0 をとる決定変数である．
\cref{eq:knapsack-objective}は選んだ品物の価値の合計を最大化することを，
\cref{eq:knapsack-capacity}は重さの合計が容量を超えないことを表す．
数理最適化の基礎については，文献~\cite{kanatani2005optimization}などを参照されたい．

%% 記号が少ない場合は，このように式の直後に「ここで，……」と本文中で定義する．
%% 記号が多い場合（大きな数理モデルなど）は，定式化の前に記号を「集合」「パラメータ」
%% 「決定変数」に分けて，description 環境で一覧にするとよい．例：
%%   \begin{description}
%%     \item[集合] $\mathcal{I} = \{1, \dots, n\}$ は品物の集合である．
%%     \item[パラメータ] $v_i$ と $w_i$ は品物 $i$ の価値と重さ，$W$ は容量である．
%%     \item[決定変数] $x_i$ は品物 $i$ を選ぶとき 1，選ばないとき 0 をとる．
%%   \end{description}

\section{既存の解法}
\label{sec:greedy}

\begin{writingnote}
既存の解法の考え方と手順を，その解法を知らない読者にもわかるように説明し，長所と短所を述べる．
短所は，後の章で述べる課題や提案手法の工夫につながるように具体的に書く．
既存の解法に自分で加えた変更は，この章ではなく「提案手法」の章に書く．
\end{writingnote}

%% アルゴリズム（擬似コード）は algorithm 環境で書く（algorithm2e パッケージ）．
%% 各行の最後には \; を付ける．
0-1 ナップサック問題の代表的な近似解法に，価値と重さの比 $v_i / w_i$ が大きい順に品物を選ぶ貪欲法がある~\cite{dantzig1957discrete}．
貪欲法の手順を\cref{alg:greedy}に示す．
貪欲法は計算が速い一方で，最適解が得られるとは限らない．

\begin{algorithm}[tbp]
  \caption{0-1 ナップサック問題に対する貪欲法}
  \label{alg:greedy}
  \KwIn{品物の価値 $v_1, \dots, v_n$，重さ $w_1, \dots, w_n$，容量 $W$}
  \KwOut{選んだ品物の集合 $S$}
  品物を $v_i / w_i$ の大きい順に並べ替える\;
  $S \gets \emptyset$，$c \gets 0$ \tcp*{$c$ は選んだ品物の重さの合計}
  \ForEach{並べ替えた順の品物 $i$}{
    \If{$c + w_i \le W$}{
      $S \gets S \cup \{i\}$\;
      $c \gets c + w_i$\;
    }
  }
  \KwRet{$S$}\;
\end{algorithm}
   % 準備（既存の知識）
%% =====================================================================
%%  第3章 関連研究
%%  書くこと：関連する研究の整理と，本研究との違い（本研究の位置づけ）
%%  章の最後で既存研究の課題を示し，次の「提案手法」につなげる．
%% =====================================================================
\chapter{関連研究}
\label{chap:related-work}

本章では，本研究に関連する研究を整理し，本研究の位置づけを明らかにする．
\begin{writingnote}
文献を 1 つずつ紹介するのではなく，いくつかの観点（問題の設定，解法の考え方など）で分類し，観点ごとに節を分けて，それぞれの考え方・長所・限界を述べるとよい．
個々の文献については，何を対象に，どのような方法で，何を明らかにしたのかを書き，本研究との共通点と相違点を述べる．
既存研究の限界は，その研究の目的や前提を踏まえて公平に，根拠を示して指摘する．
引用する文献は自分で読んで内容を確かめ，ほかの文献での紹介だけを頼りに引用すること（孫引き）は避ける．
\end{writingnote}

\section{○○に関する研究}
\label{sec:related-a}

%% 文献の引用：\cite{キー}．キーは references.bib の各文献の 1 行目に書いた名前．
%% 前の章で定義した用語や記号は，ここでそのまま使える（例：\cref{sec:problem} の問題）．
\cref{sec:problem}で述べた問題に対しては，……などの解法が研究されてきた．
例えば，Dantzig~\cite{dantzig1957discrete} は……．
また，Kellerer ら~\cite{kellerer2004knapsack} は……．

\section{△△に関する研究}
\label{sec:related-b}

%% 日本語の文献も同じように引用できる．
西~\cite{nishi2020} は……．また，西ら~\cite{nishi2011} は……．
%% 複数の文献をまとめて引用すると，番号は自動で並べ替えられ，連番は [1–3] のようにまとめられる．
このように，……に関する研究が進められている~\cite{nishi2011,nishi2020,dantzig1957discrete}．

\section{本研究の位置づけ}
\label{sec:position}

\begin{writingnote}
前の節までの整理を踏まえて，既存研究で解決されていない課題を明確に述べ，本研究がそれにどう取り組むのかを示して「提案手法」の章につなげる．
ここで示す課題は，「はじめに」の章で述べた課題と一致させる．
比較の観点は，本研究に有利なものを選ぶのではなく，問題を解くうえで重要なものを選ぶ．
「〜を扱った研究はない」と書くときは，十分に調べたうえで「筆者の知る限り」のように範囲を限って書く．
\end{writingnote}

%% 既存研究との違いは，観点を決めて表にまとめると分かりやすい．
関連研究と本研究の比較を\cref{tab:comparison}に示す．
既存研究では……が考慮されていないのに対し，本研究では……を考慮する点が異なる．

\begin{table}[tbp]
  \centering
  \caption[関連研究と本研究の比較（書式例）]{関連研究と本研究の比較（書式例．○・×は例示であり，引用文献の実際の評価ではない）}
  \label{tab:comparison}
  \begin{tabular}{lccc}
    \toprule
    & 観点 A & 観点 B & 観点 C \\
    \midrule
    文献~\cite{dantzig1957discrete} & ○ & × & × \\
    文献~\cite{nishi2011}           & ○ & ○ & × \\
    本研究                           & ○ & ○ & ○ \\
    \bottomrule
  \end{tabular}
\end{table}
    % 関連研究
%% =====================================================================
%%  第4章 提案手法
%%  書くこと：自分で考えた部分（提案）．「準備」の章の既存の内容と区別して書く．
%%  第三者が読んで実装・再現できる程度に具体的に書く．
%%  以下の手順や式は例．自分の研究に合わせて書き換える．
%% =====================================================================
\chapter{提案手法}
\label{chap:method}

本章では，○○を用いた△△問題の解法を提案する．
\begin{writingnote}
最初に提案手法の全体像（何を入力として何を出力するか，どのような手順からなるか）を示し，その後で各手順の詳細を説明する．
手順ごとに，既存の方法をそのまま用いたのか，既存の方法を変更したのか，新たに考えたのかを明記する．
既存の方法には引用を付け，引用の位置によって自分の工夫が他人の成果に見えたり，その逆になったりしないように注意する．
\end{writingnote}

\section{概要}
\label{sec:overview}

%% 最初に提案手法の全体像を図で示すと，読者が理解しやすい．
%% 簡単な図は，TikZ を使って LaTeX の中で直接描ける（下の tikzpicture 環境）．
%% 複雑な図は draw.io や PowerPoint などで作り，PDF で保存して \includegraphics で読み込むとよい．
提案手法の流れを\cref{fig:flow}に示す．
提案手法は，前処理，解の構築，解の改善の 3 つの手順からなり，解の改善は終了条件を満たすまで繰り返す．

\begin{figure}[tbp]
  \centering
  \begin{tikzpicture}[
      block/.style={draw, thick, fill=black!3, rounded corners, font=\small,
                    minimum width=36mm, minimum height=8mm, align=center},
      arrow/.style={-{Stealth[length=2.5mm]}, thick},
      node distance=5mm,
    ]
    \node[block] (input) {問題例};
    \node[block, below=of input] (pre) {前処理};
    \node[block, below=of pre] (build) {解の構築};
    \node[block, below=of build] (improve) {解の改善};
    \node[block, below=of improve] (output) {解};
    \draw[arrow] (input) -- (pre);
    \draw[arrow] (pre) -- (build);
    \draw[arrow] (build) -- (improve);
    \draw[arrow] (improve) -- (output);
    \draw[arrow] (improve.east) -- ++(8mm, 0) |- node[pos=0.25, right, font=\small, align=left]
      {終了条件を満たす\\まで繰り返す} (build.east);
  \end{tikzpicture}
  \caption{提案手法の流れ}
  \label{fig:flow}
\end{figure}

\section{各手順の詳細}
\label{sec:procedure}

\begin{writingnote}
第三者が実装し，同じ結果を再現できる程度に具体的に書く．入力と出力，各手順の処理，パラメータ，終了条件を省略しない．
設計上の選択には理由を添える．なぜその方法や評価関数を選んだのか，ほかの選択肢をなぜ採らなかったのかを説明する．
手法が前提とする仮定や，うまく働かないと考えられる場合も，「結論」の章で初めて述べるのではなく，この章で述べる．
手法の性能や実験結果は，この章ではなく「数値実験」の章に書く．
\end{writingnote}

%% 手順や記号など，並列する項目を説明するときは description 環境が便利．
各手順の内容は次のとおりである．
\begin{description}
  \item[前処理] ……により，問題例の規模を小さくする．
  \item[解の構築] \cref{alg:greedy}の貪欲法を……のように拡張した方法で，初期解を構築する．
  \item[解の改善] 構築した解を……により改善する．
\end{description}

%% 提案手法の中心となる式には番号を付け，記号は式の直後に本文中で定義する．
解の改善では，\cref{eq:penalty}の評価関数 $f$ を最大化する．
\begin{equation}
  f(\mathbf{x}) = \sum_{i=1}^{n} v_i x_i - \lambda \max\Bigl\{ 0,\ \sum_{i=1}^{n} w_i x_i - W \Bigr\}
  \label{eq:penalty}
\end{equation}
ここで，$\mathbf{x} = (x_1, \dots, x_n)$ は解を表すベクトル，
$\lambda > 0$ は容量を超えた重さに対するペナルティの重みである．
          % 提案手法（自分の提案）
%% =====================================================================
%%  第5章 数値実験
%%  書くこと：実験の目的，設定（再現に必要な情報），結果，考察
%%  以下の実験内容と数値は書き方の例（数値は架空）．自分の実験に合わせて書き換える．
%% =====================================================================
\chapter{数値実験}
\label{chap:experiments}

\begin{writingnote}
以下は記述・組版の例であり，実験設定と数値は説明用に作成した架空のものである．
図表の平均・標準偏差も実測値ではない．自分の実験内容と結果に置き換えること．
\end{writingnote}

本章では，提案手法の有効性を検証するために行った数値実験について述べる．
\begin{writingnote}
最初に，実験で何を確かめるのか（実験の目的や問い）を，「はじめに」の章で述べた目的と対応させて書く．
確かめることが複数あるときは，それぞれを番号付きの問いや実験として示し，結果と考察も同じ順に述べるとよい．
\end{writingnote}

\section{実験設定}
\label{sec:setup}

\begin{writingnote}
他の人が同じ実験を再現できるように，問題例（データ），比較手法，評価指標，パラメータ，実行環境を明記する．
問題例と比較手法は，選んだ理由も書く（例：既存研究で広く使われている問題例である，代表的な既存手法である）．
比較が公平になるように実験を計画する．確かめる点は，ソースファイルのこの位置にコメントとして書いてある．
\end{writingnote}

%% 実験を計画するときに確かめる点
%% ・調整用・検証用のデータと，最終評価用のデータを分ける．最終評価の結果で調整しない．
%% ・各手法には同じテスト問題例を用いる．計算時間・反復数などの予算と停止条件を明記する．
%% ・乱数を使う手法は複数の種で実行し，種の一覧，問題例数，問題例ごとの実行回数を書く．
%% ・問題例間のばらつきと，同じ問題例に対する乱数の種によるばらつきを区別する．
%%   平均・標準偏差がどちらを集計した値かを，表や図の説明にも明記する．
%% ・調整するパラメータの探索範囲・試行数・選択基準を記録する．

%% \paragraph{...} は目次に載らない小見出し．
\paragraph{問題例}
品物数 $n \in \{100, 200, \dots, 1000\}$ の 0-1 ナップサック問題の問題例を，$n$ ごとに 10 個ずつ生成した．
品物の価値 $v_i$ と重さ $w_i$ は 1 以上 100 以下の整数の一様乱数とし，容量 $W$ は重さの合計の半分とした．

\paragraph{比較手法}
動的計画法による厳密解法~\cite{kellerer2004knapsack}と，\cref{sec:greedy}で述べた貪欲法を比較手法とした．

\paragraph{評価指標}
最適値との差（ギャップ）と計算時間を評価指標とした．
ギャップは，最適値を $z^{*}$，得られた解の目的関数値を $z$ として，$100 (z^{*} - z) / z^{*}$ [\%] で計算した．

\paragraph{パラメータと実行環境}
%% 表は table 環境で書き，キャプションは表の「上」に付ける．
%% 罫線は \toprule・\midrule・\bottomrule（booktabs）を使い，縦罫線は使わない．
提案手法のパラメータを\cref{tab:parameters}に示す．
ペナルティの重み $\lambda$ は Optuna~\cite{akiba2019optuna} を用いて調整した．
実装には Python~\cite{python} を用い，CPU が○○，メモリが○○~GB の計算機で実験を行った\footnote{%
  実行環境の記載例：OS ○○，Python ○○，ライブラリ○○（version ○○）．}．

\begin{table}[tbp]
  \centering
  \caption{提案手法のパラメータ（架空の設定例）}
  \label{tab:parameters}
  \begin{tabular}{lc}
    \toprule
    パラメータ & 値 \\
    \midrule
    ペナルティの重み $\lambda$ & 10 \\
    解の改善の反復回数 & 1000 \\
    乱数の種 & 0 \\
    \bottomrule
  \end{tabular}
\end{table}

\section{実験結果}
\label{sec:results}

\begin{writingnote}
図表を示すだけでなく，本文で図表を参照し，そこから読み取れる全体の傾向と主な結果を文章で述べる．図表の数値をすべて文章で繰り返す必要はない．
予想と異なる結果や，提案手法が劣る結果も省かずに示す．
手法間の差を述べるときは，その差がばらつき（標準偏差など）に比べて十分に大きいかを確かめる．
結果の解釈（なぜその結果になったのか）は，次の「考察」の節で述べる．
\end{writingnote}

%% 図は figure 環境で書き，キャプションは図の「下」に付ける．図のファイルは figures/ に置く．
%% グラフは PDF（ベクター形式）で保存すると，拡大しても粗くならない．
%% 図を並べるときは subfigure 環境を使う．\cref で「図5.1(a)」のように参照できる．
各手法のギャップと計算時間を\cref{fig:results}に示す．
\cref{fig:results-gap}より，提案手法のギャップはすべての $n$ で貪欲法より小さい．
また，\cref{fig:results-time}より，厳密解法の計算時間は $n$ とともに急激に増加するのに対し，
提案手法の計算時間はほぼ $n$ に比例して緩やかに増加する．

\begin{figure}[htbp]
  \centering
  \begin{subfigure}{0.48\linewidth}
    \centering
    \includegraphics[width=\linewidth]{results-gap.pdf}
    \caption{最適値とのギャップ}
    \label{fig:results-gap}
  \end{subfigure}
  \hfill
  \begin{subfigure}{0.48\linewidth}
    \centering
    \includegraphics[width=\linewidth]{results-time.pdf}
    \caption{計算時間}
    \label{fig:results-time}
  \end{subfigure}
  \caption[品物数と各手法の性能（架空データ）]{品物数 $n$ と各手法の性能（架空データ：各 $n$ で 10 個の問題例の平均を想定）}
  \label{fig:results}
\end{figure}

$n = 1000$ のときの結果を\cref{tab:results}に示す．各列の最良値を太字で示している．

\begin{table}[htbp]
  \centering
  \caption[$n = 1000$ における結果（架空データ）]{$n = 1000$ における各手法の結果（架空データ：10 個の問題例の平均 $\pm$ 標準偏差を想定）}
  \label{tab:results}
  \begin{tabular}{lrr}
    \toprule
    手法 & ギャップ [\%] & 計算時間 [s] \\
    \midrule
    厳密解法 & $\mathbf{0.00 \pm 0.00}$ & $83.0 \pm 21.7$ \\
    貪欲法   & $1.07 \pm 0.35$ & $\mathbf{0.01 \pm 0.00}$ \\
    提案手法 & $0.22 \pm 0.08$ & $2.00 \pm 0.12$ \\
    \bottomrule
  \end{tabular}
\end{table}

\section{考察}
\label{sec:discussion}

\begin{writingnote}
結果の解釈を述べる．結果を繰り返すだけにならないように注意する．
主な結果ごとに，なぜその結果になったのかを説明する．推測にとどまる説明は「〜と考えられる」のように書き，実験で確かめた事実と区別する．
既存研究で報告されている結果と比べられる場合は，一致する点と異なる点を述べる．
そのうえで，どのような条件で提案手法が有効なのか，どのような場合には有効でないのか（限界）を，結果を根拠に議論する．
ここで述べた限界は，「結論」の章の今後の課題につながる．
\end{writingnote}

本実験の条件では，……について……という傾向が見られた．
一方，……では……であり，適用には……という制約がある．
     % 数値実験
%% =====================================================================
%%  第6章 結論
%%  書くこと：本研究のまとめ（目的・方法・結果）と今後の課題
%%  本文に書いていない新しい結果をここで初めて出さない．
%% =====================================================================
\chapter{結論}
\label{chap:conclusion}

\section{まとめ}
\label{sec:summary}

\begin{writingnote}
「はじめに」の章で述べた目的と課題を振り返り，それに対して何を行い，何がわかったかを要約する．目的と課題は「はじめに」の章と同じ言葉で書くと，読者が対応を追いやすい．
主な結果は条件や数値を挙げて具体的に述べ，目的がどこまで達成できたかを明確にする．
そのうえで，「関連研究」の章で整理した既存研究に対して，本研究が何を新たに加えたのかを述べる．
数値実験で確かめた条件を超えて結論を一般化しない．
\end{writingnote}

本研究では，……を目的として，……を提案した．
数値実験の結果，提案手法は……であることを確認した．

\section{今後の課題}
\label{sec:future-work}

\begin{writingnote}
今後の課題は，本研究の限界のうち，今後の研究で取り組めるものから導く．限界は「提案手法」や「数値実験」の章で述べたものを取り上げ，この節で初めて持ち出さない．
「さらなる改良が必要である」のような漠然とした書き方ではなく，何が未解決で，どのように取り組むかを具体的に書く．
\end{writingnote}

%% まだ書いていない所には \TODO{...} を置いておくとよい．
%% 例：\TODO{今後の課題を具体的に書く}
%% review オプションを付けると赤字で表示され，付けないと表示されずに警告が出る．
今後の課題として，……が挙げられる．
      % 結論

%% =====================================================================
%%  謝辞：研究でお世話になった方々への感謝を述べる．
%% =====================================================================
\begin{acknowledgments}
本研究を進めるにあたり，終始丁寧なご指導をいただいた○○ ○○先生に深く感謝いたします．
また，日頃から議論に付き合っていただいた知能システム研究室の皆様に感謝いたします．
\end{acknowledgments}
   % 謝辞
\bibliography{references}          % 参考文献（references.bib から自動で作られる）
%% =====================================================================
%%  研究業績
%%  在学中の論文や学会発表を，種類ごとに新しい順に書く．
%%  業績がない場合は，main.tex の %% =====================================================================
%%  研究業績
%%  在学中の論文や学会発表を，種類ごとに新しい順に書く．
%%  業績がない場合は，main.tex の %% =====================================================================
%%  研究業績
%%  在学中の論文や学会発表を，種類ごとに新しい順に書く．
%%  業績がない場合は，main.tex の \input{chapters/publications} の行を削除する．
%%  各項目には，著者（自分の名前に下線），題目，論文誌名または会議名，
%%  巻・号・ページ，開催地，年月を書く．書き方の細部は指導教員の指示に従う．
%% =====================================================================
\begin{publications}

\section*{学術論文}
\begin{enumerate}
  \item \underline{法政 太郎}，劉 子昂：○○を用いた△△問題の解法，
    ○○学会論文誌，Vol.~○，No.~○，pp.~○--○，2027．
\end{enumerate}

\section*{国際会議}
\begin{enumerate}
  \item \underline{T.~Hosei} and Z.~Liu: ``Title of the paper,''
    in \textit{Proceedings of the ○○ International Conference on ○○}, pp.~○--○, ○○, 2026.
\end{enumerate}

\section*{国内会議・研究会}
\begin{enumerate}
  \item \underline{法政 太郎}，劉 子昂：○○を用いた△△問題の解法，
    ○○学会 ○○大会，○○，2026 年○月．
\end{enumerate}

\end{publications}
 の行を削除する．
%%  各項目には，著者（自分の名前に下線），題目，論文誌名または会議名，
%%  巻・号・ページ，開催地，年月を書く．書き方の細部は指導教員の指示に従う．
%% =====================================================================
\begin{publications}

\section*{学術論文}
\begin{enumerate}
  \item \underline{法政 太郎}，劉 子昂：○○を用いた△△問題の解法，
    ○○学会論文誌，Vol.~○，No.~○，pp.~○--○，2027．
\end{enumerate}

\section*{国際会議}
\begin{enumerate}
  \item \underline{T.~Hosei} and Z.~Liu: ``Title of the paper,''
    in \textit{Proceedings of the ○○ International Conference on ○○}, pp.~○--○, ○○, 2026.
\end{enumerate}

\section*{国内会議・研究会}
\begin{enumerate}
  \item \underline{法政 太郎}，劉 子昂：○○を用いた△△問題の解法，
    ○○学会 ○○大会，○○，2026 年○月．
\end{enumerate}

\end{publications}
 の行を削除する．
%%  各項目には，著者（自分の名前に下線），題目，論文誌名または会議名，
%%  巻・号・ページ，開催地，年月を書く．書き方の細部は指導教員の指示に従う．
%% =====================================================================
\begin{publications}

\section*{学術論文}
\begin{enumerate}
  \item \underline{法政 太郎}，劉 子昂：○○を用いた△△問題の解法，
    ○○学会論文誌，Vol.~○，No.~○，pp.~○--○，2027．
\end{enumerate}

\section*{国際会議}
\begin{enumerate}
  \item \underline{T.~Hosei} and Z.~Liu: ``Title of the paper,''
    in \textit{Proceedings of the ○○ International Conference on ○○}, pp.~○--○, ○○, 2026.
\end{enumerate}

\section*{国内会議・研究会}
\begin{enumerate}
  \item \underline{法政 太郎}，劉 子昂：○○を用いた△△問題の解法，
    ○○学会 ○○大会，○○，2026 年○月．
\end{enumerate}

\end{publications}
      % 研究業績（業績がない場合はこの行を削除する）

\appendix                          % ここから付録（付録A，付録B，...）
%% =====================================================================
%%  付録
%%  本文に入れると流れを妨げる詳細（証明，追加の実験結果，ソースコードなど）を書く．
%%  \appendix の後の \chapter は「付録A」「付録B」… となる．
%% =====================================================================
\chapter{ソースコード}
\label{chap:source-code}

%% ソースコードは lstlisting 環境で書く．language= で言語を指定すると色が付く．
%% ファイルから読み込むこともできる：\lstinputlisting[language=Python]{code/main.py}
\cref{alg:greedy}の貪欲法を Python で実装した例を\cref{lst:greedy}に示す．

\begin{lstlisting}[language=Python, caption={貪欲法の実装例}, label=lst:greedy]
def greedy_knapsack(values, weights, capacity):
    """価値と重さの比が大きい順に品物を選ぶ"""
    order = sorted(range(len(values)),
                   key=lambda i: values[i] / weights[i], reverse=True)
    selected, total_weight = [], 0
    for i in order:
        if total_weight + weights[i] <= capacity:  # 容量を超えなければ選ぶ
            selected.append(i)
            total_weight += weights[i]
    return selected
\end{lstlisting}

\chapter{記号一覧}
\label{chap:notation}

本論文で用いる主な記号を\cref{tab:notation}に示す．
スカラーは細字の斜体（例：$x$），ベクトルは太字の小文字（例：$\mathbf{x}$），
行列は太字の大文字（例：$\mathbf{A}$），集合は花文字（例：$\mathcal{I}$）で表す．
確率変数を用いる場合は，確率変数を大文字（例：$X$），その値を小文字（例：$x$）で表す．

%% 記号の書き方は，分野の標準的な教科書に合わせるとよい
%% （例：強化学習なら Sutton と Barto の教科書）．
\begin{table}[htbp]
  \centering
  \caption{主な記号}
  \label{tab:notation}
  \begin{tabular}{ll}
    \toprule
    記号 & 意味 \\
    \midrule
    $n$ & 品物の数 \\
    $v_i$, $w_i$ & 品物 $i$ の価値，重さ \\
    $W$ & 容量 \\
    $x_i$ & 品物 $i$ を選ぶとき 1，選ばないとき 0 をとる決定変数 \\
    $\mathbf{x}$ & 解を表すベクトル $(x_1, \dots, x_n)$ \\
    $f(\mathbf{x})$ & 提案手法の評価関数 \\
    $\lambda$ & ペナルティの重み \\
    \bottomrule
  \end{tabular}
\end{table}


\end{document}
